\documentclass[11pt]{article}
\usepackage[T1]{fontenc}
\usepackage[utf8]{inputenc}
\usepackage{lmodern,amsmath,amssymb,amsthm}
\usepackage[margin=1in]{geometry}
\usepackage[hidelinks]{hyperref}
\setlength{\emergencystretch}{1em}
\newtheorem{theorem}{Theorem}
\newtheorem{lemma}[theorem]{Lemma}
\newcommand{\Z}{\mathbb Z}\newcommand{\R}{\mathbb R}
\newcommand{\C}{\mathbb C}\newcommand{\T}{\mathbb T}
\newcommand{\E}{\mathbb E}\newcommand{\poly}{\operatorname{poly}}
\newcommand{\Lip}{\operatorname{Lip}}
\title{Equal-scale factorization and the bias and torsion joins}
\author{Author proposal: /root/astra\_v14\_baseline\_compiler}
\date{Protocol 1.4; separate successor to the one-variable D/F packet}
\begin{document}\maketitle
\noindent This packet does not revise the frozen D/F or Section 13 proofs.
It supplies a proposed connecting proof for the previously open
D.5-to-E.11/E.13 dependency. All acceptance and incorporation decisions
remain independent. The only quantitative equidistribution input is the
\emph{final published} equal-side theorem in [GT14, pp.1176--1177], with
the published minor corrections to [GT12]. No unpublished arbitrary-side
extension is imported.

\section{Conventions and exact statements}
Fix standard positive integers $k,d$ and an unbounded integer $N$.
All groups, rational structures, lattices, coordinate bases and finite
dimensions are standard, with natural ultrapower extensions. A nilpotent
Lie group is connected and simply connected. The rational degree filtration
$G_\bullet$ has finite length; the source $\Z^k$ has its total-degree
filtration. Multiindices have nonnegative integer entries, and
$|\alpha|=\sum_i\alpha_i$. All vertical identities for a limit nilcharacter
hold at every limit group element and point, as in the accepted
tautification-003 convention. We do not substitute equivariance only at
standard points.

In fixed Mal'cev coordinates, small polynomial coefficients mean bounded
constant coefficient and coefficient $O(N^{-|\alpha|})$ at positive index
$\alpha$ in either the monomial or product-binomial basis. The constants
are standard and may depend on the actual fixed geometric data and orbit.
Total equidistribution on $[N]^k$ means equidistribution on every product
of dense arithmetic progressions in $[N]$, with standard Lipschitz tests.
It also applies to internal tests with a common standard Lipschitz bound:
finite uniform nets in a compact Lipschitz ball reduce those tests to
standard ones, with the desired accuracy chosen before taking the limit.

\begin{theorem}[Equal-scale factorization, MF-D4]\label{thm:factor}
For $g\in{}^*\poly(\Z^k,G_\bullet)$ there are a standard connected rational
subgroup $P\leq G$ and internal polynomial maps $\epsilon,p,\gamma$ such that
\[
                       g(n)=\epsilon(n)p(n)\gamma(n)
                         \qquad(n\in{}^*\Z^k).                 \tag{1}
\]
The map $p$ lies in $P$ with its induced filtration, and is totally
equidistributed on $[N]^k$ in $P/\Delta$ for every standard finite-index
$\Delta\leq P\cap\Gamma$, with the same $p$ and $P$.
The coefficients of $\epsilon$ are small as specified above; it is bounded
and $O(1/N)$-Lipschitz on every fixed standard dilate $[-AN,AN]^k$.
Its group Taylor coefficient at positive index $\alpha$ is
$O(N^{-|\alpha|})$ as well. There is a standard discrete subgroup
$\Lambda\geq\Gamma$ containing all values of $\gamma$, and a standard
$q>0$ such that
\[
           \gamma(n+qe_i)\Gamma=\gamma(n)\Gamma
                  \qquad(n\in{}^*\Z^k, 1\leq i\leq k).       \tag{2}
\]
There is no properness requirement on $P$, and no claim that $\gamma$ is
$\Gamma$-valued. All factors in (1) are defined on the full integer lattice.
\end{theorem}

Write $I=[-N,N]\cap{}^*\Z$.
\begin{theorem}[Actual bias and torsion outputs, MF-E11/E13]\label{thm:outputs}
Let $d\geq2$ be standard, and let $\chi$ be a globally represented unit
vector nilcharacter of degree $d$ on ${}^*\Z^k$.
If a scalar coordinate of $\chi$ has nonzero correlation with a degree
at most $d-1$ nilsequence after multiplication by the indicator of a
convex polytope in $I^k$ and a product of dense progressions in $[N]$,
then $\chi$ agrees on all of $I^k$ with a globally represented
degree-at-most-$(d-1)$ vector nilsequence. The same conclusion holds if
$\chi^{\otimes a}$ agrees on $I^k$ with a degree-at-most-$(d-1)$ vector
nilsequence for a standard integer $a\geq1$.

These assertions give exactly the lower-degree conclusions of [GTZ, E.11]
and [GTZ, E.13] on their required symmetric box. They include finite
scalar-coordinate reductions for vector correlations. They do not require
an extrapolation from an unrelated subprogression or the deletion of a
rational Taylor coefficient.
\end{theorem}

\section{A weighted inverse-affine estimate}
The following estimate also supplies the large-offset transition behind
[GT12, Proposition 9.2]. Merely inserting an offset of size $N$ into the
height bound of its Lemma 8.4 would not give the required standard bound.

\begin{lemma}[Large offsets, bounded steps, MF-AFFINE]\label{lem:affine}
Let $f(n)$ be a real polynomial of fixed total degree at most $D$ in $k$
variables. Let $A_i$ be internal integers with $|A_i|=O(N)$ and let $b_i$
be nonzero integers of bounded absolute value. Put
$\widetilde f(m)=f(A_1+b_1m_1,\ldots,A_k+b_km_k)$.
Suppose every positive product-binomial coefficient $\widetilde f_\beta$
satisfies $\|\widetilde f_\beta\|_\T=O(N^{-|\beta|})$.
There is a standard positive integer $r$, depending only on $D,k$ and the
step bounds, such that every positive coefficient of $rf$ has the same
bound $O(N^{-|\alpha|})$ modulo one.
\end{lemma}
\begin{proof}
Invert the substitution over $\R$ and expand
\[
 \prod_i\binom{(n_i-A_i)/b_i}{\beta_i}
        =\sum_{\alpha\leq\beta}c_{\alpha\beta}\binom n\alpha.
                                                               \tag{3}
\]
Expanding each falling factorial first in monomials shows that the
denominators divide a fixed integer formed from $\beta_i!$ and
$b_i^{\beta_i}$. Conversion of ordinary powers to binomial coefficients
uses integer Stirling coefficients. Thus one standard $r$ clears all
denominators, independently of the integer numerators $A_i$.
The coefficient of $n^\nu$ before the last conversion is a polynomial in
the $A_i$ of total degree at most $|\beta|-|\nu|$. Consequently
\[
 |c_{\alpha\beta}|=O(N^{|\beta|-|\alpha|}),
 \qquad rc_{\alpha\beta}\in{}^*\Z.                             \tag{4}
\]
The case $\beta=0$ contributes only to $\alpha=0$.
For each $\beta>0$ write $\widetilde f_\beta=z_\beta+e_\beta$, where
$z_\beta$ is an integer and $e_\beta=O(N^{-|\beta|})$.
In the coefficient of index $\alpha>0$ in $rf$, the contributions
$rc_{\alpha\beta}z_\beta$ are integers, while
\[
 \sum_{\beta\geq\alpha}rc_{\alpha\beta}e_\beta
                         =O(N^{-|\alpha|}).                    \tag{5}
\]
There are only a fixed number of indices. This proves the claim. Negative
steps cause no change after taking a positive common denominator. The
uncontrolled constant coefficient causes no positive-index error.
\end{proof}

\section{From a dense rectangle to the published equal-side theorem}
\begin{lemma}[Weak factorization in finitely many parameters]\label{lem:weak}
Suppose $G_0=G_1=G$. If $g\Gamma$ is not totally equidistributed on $[N]^k$,
it has a factorization (1) with $p$ in a standard proper connected rational
subgroup, retaining all conclusions of Theorem~\ref{thm:factor} except
the final equidistribution condition.
\end{lemma}
\begin{proof}
Failure supplies a product of progressions
$P_i=\{a_i+b_it:1\leq t\leq L_i\}\subseteq[N]$ with $L_i\geq\delta N$,
standard $\delta>0$, and a bounded Lipschitz test with noninfinitesimal
discrepancy. Their integer steps obey $1\leq b_i\leq2/\delta$ and their
offsets are $O(N)$. Subtract the Haar mean, normalize the test and rotate a
constant complex phase so that the real mean discrepancy is positive.

In the index rectangle $\prod[ L_i]$, partition each axis into blocks of a
common length $L=\lfloor\eta N\rfloor$, leaving a final remainder of length
less than $L$. Choose the standard $\eta>0$ small compared with the
discrepancy, $\delta$, $k$ and the test norm. The discarded proportion is
at most $\sum_iL/L_i\leq k\eta/\delta+o(1)$, so its contribution is less
than one quarter of that discrepancy. At least one full cube of side $L$
has a positive standard discrepancy. On this cube the polynomial is
$\widetilde g(m)=g(A_1+b_1m_1,\ldots,A_k+b_km_k)$ with $|A_i|=O(N)$.

Apply the exact equal-side theorem of [GT14, pp.1176--1177] at length $L$,
using a standard error parameter below this discrepancy and below the
reciprocal rationality bound of the fixed adapted basis. The result gives
a nonzero bounded horizontal integer character $\xi$ for which all positive
coefficients of the real horizontal lift of $\xi\widetilde g$ have circle
norm $O(L^{-|\alpha|})=O(N^{-|\alpha|})$.
Only finitely many such characters are possible, so the selected character
is standard. Lemma~\ref{lem:affine} replaces it by a nonzero standard
multiple and gives the same coefficient bounds for $\xi g$ on the
original variables. This is the precise rectangle-to-equal-cube join;
the original unequal-side theorem is not used.

Normalize $g(0)=c\lambda$ with bounded $c$ and $\lambda\in{}^*\Gamma$
and apply the preceding argument to $c^{-1}g\lambda^{-1}$, whose failure
is preserved by bounded left translation. Its coordinates are
$\sum_{\alpha>0}t_\alpha\binom n\alpha$, with $t_\alpha$ supported in
$G_{|\alpha|}$ by [GT12, Lemma 6.7]. Write the real horizontal lift as
$u\cdot\psi$, with $u$ a fixed integer vector.
Choose integers $a_\alpha$ nearest to $u\cdot t_\alpha$.
On a support where $u$ is nonzero choose a fixed rational vector
$v_\alpha$ of that support with $u\cdot v_\alpha=1$, and define
\[
 w_\alpha=(u\cdot t_\alpha-a_\alpha)v_\alpha,
 \quad z_\alpha=a_\alpha v_\alpha,
 \quad \psi(\epsilon_0(n))=\sum w_\alpha\binom n\alpha,
 \quad \psi(\gamma_0(n))=\sum z_\alpha\binom n\alpha.           \tag{6}
\]
If $u$ vanishes on a support, set all these vectors and its integer to zero.
Coordinate support and polynomial group closure make these maps polynomial
for the original filtration. The real lift is identically zero on
$p=\epsilon_0^{-1}c^{-1}g\lambda^{-1}\gamma_0^{-1}$.
Thus $p$ lies in its standard proper connected rational kernel.
The original factors are $\epsilon=c\epsilon_0$ and
$\gamma=\gamma_0\lambda$. The $w_\alpha$ have precisely the required
small bounds; the $z_\alpha$ have a common standard denominator, without
a numerator bound. Controlled products and periods are checked next.
\end{proof}

\section{Controlled products, periods, and descent}
Scale all variables by $n_i=Nt_i$. Fixed coordinate polynomials with small
coefficients become polynomials in $t$ with bounded coefficients. The
fixed BCH and coordinate-change polynomials preserve this property under
finite products, inverses and fixed rational subgroup embeddings.
Conversion between monomial and product-binomial bases is triangular in
componentwise indices, so the coefficient of $\binom n\alpha$ is
$O(N^{-|\alpha|})$. Differentiating the scaled coordinate maps on the
compact box $[-A,A]^k$ gives the claimed $O(1/N)$ Lipschitz bound on every
fixed dilate. A real-coordinate polynomial extension of $\epsilon$ has the
same bounds, which will be used below.

These estimates also imply the group Taylor bounds. Assign formal weight
$|\alpha|$ to each logarithmic coefficient of $n^\alpha$ and weight zero to
the bounded constant term. Group Taylor extraction in any fixed admissible
order uses fixed BCH polynomials. Inductively every term of formal weight
$w$ has coordinate polynomial total degree at most $w$; removing a Taylor
factor at index $\beta$ preserves this property because its binomial power
has degree $|\beta|$. The coefficient at index $\alpha$ cannot contain a
term of formal weight less than $|\alpha|$: such a term has total degree
less than $|\alpha|$ and has zero coefficient at that index. This is the
same triangular Taylor extraction as [GTZ, B.9], now with product-binomial
indices. Each surviving BCH monomial contributes at least
$N^{-|\alpha|}$ decay. In particular no estimate is inferred solely from
separate degree in individual variables.

The common discrete overgroup argument is independent of the number of
parameters. In a rational lower-central basis of weights $w_i$, choose a
standard $D$ clearing every coordinate denominator and every fixed BCH
coefficient. As in DF002, the additive lattice
$\sum_iD^{-(2w_i-1)}\Z X_i$ is BCH-closed: a nonlinear term with $t$
inputs has output denominator exponent exceeding the input sum by at least
$t-1$. Its exponential is a standard discrete overgroup $\Lambda\geq\Gamma$
containing the rational factors and any finite product of them.

Here is the full multivariable period argument, in the right-coset direction.
Each coordinate of $\gamma(n)$ is a polynomial of fixed degree whose values
have a common standard denominator. Tensor-product interpolation on a fixed
integer grid gives a standard common denominator for all its coefficients.
For every $i$ the vector
\[
 R_i(n,r)=\psi\bigl(\gamma(n)^{-1}\gamma(n+re_i)\bigr)
\]
is a rational coefficient polynomial with a common standard denominator.
It is identically zero when $r=0$. Choose a standard positive $q$ divisible
by all these denominators. Each monomial then has a positive power of $r$,
so $R_i(n,q)$ is integral for every internal integer vector $n$.
This is exactly (2). This proof does not use the wrongly oriented quotient
in the printed proof of [GT12, A.12]. The corresponding one-variable
sublemma of DF002 has a separate independent root review.

To finish Theorem~\ref{thm:factor}, first normalize the constant term and
restrict to $G_1$ if $G_1\ne G$. The identity
\[
 c^{-1}g(n)\lambda^{-1}
          =\lambda\bigl(g(0)^{-1}g(n)\bigr)\lambda^{-1}\in G_1
\]
holds in any number of variables, since first derivatives lie in $G_1$ and
the positive groups are normal in $G_0$. Thus a controlled representation
exists with a taut positive filtration.
Choose a standard rational representing subgroup of least dimension.
For any standard finite-index lattice $\Delta\leq P\cap\Gamma$, failure
of total equidistribution in $P/\Delta$ permits Lemma~\ref{lem:weak}, after
the same normalization if its induced first group is proper.
It supplies $p=\epsilon_2p_2\gamma_2$ in a proper rational $P_2<P$.
The correct new identity is
\[
                 g=(\epsilon\epsilon_2)p_2(\gamma_2\gamma).
\]
The preceding controlled-product and discrete-overgroup arguments verify
every required bound and period. Rationality relative to $\Delta$ is the
same rational structure as relative to $P\cap\Gamma$.
This contradicts minimal dimension. It proves total equidistribution for
every standard finite-index cover with the same factors. Only a minimum
among $0,\ldots,\dim G$ was selected; no infinite sequence of refinements
or dependence of the chosen factors on the later cover is hidden here.

\section{One quotient for all periodic residue classes}
We now work in a taut ambient degree-$d$ group $G$, so $G_d$ is central.
Tautification-003, Theorem 1, supplies this exact representation for each
nilsequence used below with the same positive-degree groups and the same
standard vertical character. Products of these representations remain taut.

Apply Theorem~\ref{thm:factor}, and choose a standard representative
$\lambda_\rho\in\Lambda$ for each value of the right coset of $\gamma(n)$
on a residue class $n\equiv\rho\pmod q$, $\rho\in(\Z/q\Z)^k$.
Such representatives can be standard: $\Lambda/\Gamma$ is a fixed finite
set, and periodicity makes the coset constant on each class. Let
\[
 \Gamma_0=\bigcap_{\lambda\in\Lambda}\lambda\Gamma\lambda^{-1},
 \qquad \Delta=P\cap\Gamma_0.                                \tag{7}
\]
The first intersection is the kernel of the action on the finite coset set,
so it is a standard normal finite-index subgroup of $\Lambda$ contained in
$\Gamma$. Thus $\Delta$ is finite index in $P\cap\Gamma$, and $p$ is
totally equidistributed on $P/\Delta$. For every $\rho$ the function
\[
                      x\longmapsto F(zx\lambda_\rho\Gamma)   \tag{8}
\]
is well-defined on this \emph{same} quotient: for $\delta\in\Delta$,
$\lambda_\rho^{-1}\delta\lambda_\rho\in\Gamma$.
No deletion of a rational horizontal coefficient is involved.

\begin{lemma}[A nonzero integral from the actual cutoff]\label{lem:integral}
Let $F$ be an internal bounded Lipschitz scalar test of standard norm on
$G/\Gamma$. If
\[
 \left|\E_{n\in I^k}1_\Omega(n)1_{P_1\times\cdots\times P_k}(n)
                           F(g(n)\Gamma)\right|>c>0           \tag{9}
\]
for a convex polytope $\Omega\subseteq I^k$ and dense progressions
$P_i\subseteq[N]$, then for some residue $\rho$ and bounded
$z=\epsilon(n_0)$ one has
\[
          \left|\int_{P/\Delta}F(zx\lambda_\rho\Gamma)
                                           \,d\mu(x)\right|>c'>0. \tag{10}
\]
Here $c'$ is standard and positive. The factors and quotient are those
just constructed for the same actual $g$.
\end{lemma}
\begin{proof}
Partition the progression product into its finitely many nonempty residue
classes modulo $q$. Each intersection is again a product of dense
progressions, by the elementary congruence calculation for each axis.
Their steps are bounded by the old step times $q$. Rotate a common complex
phase in (9); one class retains a real sum at least $c_0N^k$ for a
standard $c_0>0$. Fix this lower bound before selecting a mesh.

Use index coordinates on that progression product. The inverse image of
$\Omega$ intersected with the index rectangle is convex and lies in a box
of side $O(N)$. Partition it by cubes of side $\lfloor\eta N\rfloor$.
The cubes meeting its boundary have at most $O_k(\eta N^k)+O_k(N^{k-1})$
integer points, with constants also depending on the fixed step bounds.
Indeed place disjoint unit cubes at those integer points and enclose them
in the $O_k(\eta N+1)$ neighborhood of the convex boundary; apply the
published boundary-volume estimate [GT10, Corollary A.2, p.1810].
For a convex set with empty interior, first thicken by an arbitrarily small
ball and then take the decreasing limit of its boundary neighborhoods.
Choose standard $\eta$ small enough that the discarded boundary
contribution is below $c_0N^k/4$. The total number of points in the full
cubes is at most $N^k$. Some cube contained in the convex set therefore
gives a dense progression product $R$ whose phase-adjusted average is at
least $3c_0/4$.

The choice of $\eta$ can also ensure that
$|F(\epsilon(n)x\lambda_\rho\Gamma)-F(\epsilon(n_0)x\lambda_\rho\Gamma)|$
is below $c_0/4$ for all $n,n_0\in R$ and all
$x\in P/\Delta$. This follows from the $O(1/N)$ variation of $\epsilon$,
its bounded image, the fixed standard representative $\lambda_\rho$, and
compactness of $P/\Delta$. Choose the mesh only after these standard
constants and the original positive discrepancy are known.
For $n\in R$, (1) and the fixed right coset give
$F(g(n)\Gamma)=F(\epsilon(n)p(n)\lambda_\rho\Gamma)$.
Freeze $\epsilon$ at $n_0$. Total equidistribution of the same $p$ in
$P/\Delta$ on $R$ identifies the resulting average with the integral in
(10), up to an infinitesimal error. Internal-test uniformity follows from
the finite Lipschitz net specified at the start. This proves (10).
\end{proof}

\section{Exact reconstruction on the symmetric box}
Let $G=G^{\chi}\times G^0$ be the taut product of a degree-$d$ nilcharacter
representation and any degree-at-most-$(d-1)$ auxiliary representation.
The latter has trivial degree-$d$ group. Write $F_\chi$ for the original
unit vector test, and let $\eta$ be its real standard vertical character.
Set
\[
                K=P_d=P\cap G_d.
\]
We justify the rational intersection with the original lattice explicitly.
Fix a rational Mal'cev basis for $G/\Gamma$. A connected subgroup $H$
with compact $H/(H\cap\Gamma)$ has rational Lie algebra in this basis.
Indeed, [GT12, Lemma A.2] makes the coordinates of $\log\gamma$ rational
for $\gamma\in\Gamma$. The Lie algebra $\mathfrak l$ generated by
$\log(H\cap\Gamma)$ therefore has a rational basis, by selecting finitely
many independent iterated brackets. If $\mathfrak l$ were proper in
$\mathfrak h$, put it in a maximal proper subalgebra $\mathfrak m$.
The last lower central term not contained in $\mathfrak m$ supplies a
vector outside $\mathfrak m$ normalizing it. Thus maximality makes
$\mathfrak m$ an ideal, and $\mathfrak h/\mathfrak m$ has dimension one,
since otherwise a line would be a proper nonzero subalgebra. The quotient
functional exponentiates by BCH to a surjection $H\to\mathbb R$ killing
$H\cap\Gamma$, contradicting compactness. Hence $\mathfrak l=\mathfrak h$.
Apply this argument to $P$ and $G_d$. Their Lie-algebra intersection has
a rational basis by rational linear algebra. The global exponential
homeomorphism and BCH identify $K$ with its exponential, so $K$ is closed
and connected. Equip $K$ with its rational lower central filtration.
Choose a standard $Q$ bounding the finitely many rational coordinates.
[GT12, Definition 2.5 and Proposition A.10] then give a Mal'cev basis for
$K/(K\cap\Gamma)$, so this actual intersection is a lattice. The trivial
subgroup case is immediate. Since $G_d$ is central in the taut ambient
group, $K$ is central too.
The coset map from $(K\cap\Gamma)/(K\cap\Delta)$ to
$(P\cap\Gamma)/\Delta$ is injective. Thus $K\cap\Delta$ has finite index
in $K\cap\Gamma$ and is a lattice in $K$, without requiring normality
of $\Delta$. The central-quotient argument in LP13-002, Section 2, now
applies with its rationality premise established: the image
$\overline\Delta$ of $\Delta$ is a lattice in $\overline P=P/K$.
The quotient filtration kills $P_d=K$ and hence has degree at most $d-1$.
All these objects depend only on the fixed standard geometry.

Suppose that
\[
                  \eta(\operatorname{pr}_\chi K)=0.            \tag{11}
\]
For each $\rho$ and bounded real-coordinate parameter $t\in[-1/5,1/5]^k$,
define, initially on $P/\Delta$,
\[
 H_\rho(x,t)=F_\chi\bigl(
   \widetilde\epsilon_\chi(10Nt)\,x_\chi\,(\lambda_\rho)_\chi
                                                    \Gamma_\chi\bigr), \tag{12}
\]
where $\widetilde\epsilon$ is the real-coordinate polynomial extension.
This is well-defined by (7). It is invariant under $K$ by the all-limit
vertical identity and (11); centrality permits its element to pass across
$\widetilde\epsilon$ and $\lambda_\rho$. Hence (12) descends to
$\overline P/\overline\Delta$.
The scaled-coefficient estimates imply a standard Lipschitz bound in $t$,
uniformly in $N$. Smooth multiplication on fixed compact ranges gives the
same conclusion in $x$. Invariance descends this bound to the quotient
metric; any fixed standard Riemannian metric is equivalent on this compact
nilmanifold.

Choose a fixed Lipschitz function $B(t)$ on $\T^k$ equal to one on
$[-1/10,1/10]^k$ and zero outside $(-1/5,1/5)^k$.
Using centered representatives inside that chart, define
$\widetilde H_\rho(x,t)=B(t)H_\rho(x,t)$ there and zero outside.
It is a globally bounded Lipschitz vector function. The zero collar makes
the definition independent of the chart boundary. The functions are limit
functions with a standard bound, since their coordinate-level construction
has uniform bounds before taking the ultrapower.

For each $\rho\in(\Z/q\Z)^k$ choose a standard Lipschitz function
$b_\rho$ on $\T^k$ with
$b_\rho(\sigma/q)=1_{\rho=\sigma}$. For example use products of triangular
bumps of radius $1/q$; for $q=1$ use the constant function one.
The vector function
\[
 \mathcal H(x,t,z)=\sum_\rho b_\rho(z)\widetilde H_\rho(x,t)    \tag{13}
\]
is globally bounded Lipschitz with a standard bound. Its orbit is
\[
          n\longmapsto\bigl(\overline p(n)\overline\Delta,
                                      n/(10N),\,n/q\bigr).    \tag{14}
\]
This orbit has degree at most $d-1$: its first component has that degree
and both tori have degree one, with $1\leq d-1$.
For every $n\in I^k$, exactly the correct residue value is sampled in
(13), and $B(n/(10N))=1$. The identities (1), (2), and (12) give
\[
                        \chi(n)=\mathcal H((14)).              \tag{15}
\]
This is an equality of the actual vector functions on the entire symmetric
box. It is not an extrapolation of an equality on just one progression.
No unit condition is imposed on $\mathcal H$ away from the orbit; a
lower-degree vector nilsequence does not require one.

\section{Bias and torsion applications on the same objects}
\begin{proof}[Proof of Theorem~\ref{thm:outputs}]
For the bias assertion, choose the scalar coordinate and the lower-degree
scalar test occurring in the actual nonzero correlation. If a vector
correlation is given, a fixed finite coordinate choice supplies one such
scalar nonzero average; its density loss is standard. Form the taut product
representation $G^\chi\times G^0$ and the product test
$F= (F_\chi)_i\overline{F_0}$, using the sign in that correlation.
It has vertical character $\eta$ on the degree-$d$ group because $G^0_d$
is trivial. Apply Theorem~\ref{thm:factor} and Lemma~\ref{lem:integral} to
this actual product orbit and test, obtaining a nonzero integral (10).

Translation of Haar measure on $P/\Delta$ by any $u\in K$ preserves that
integral. The integrand is multiplied by
$e(\eta(\operatorname{pr}_\chi u))$, since $K$ is central in the ambient
taut group. The integral is nonzero, so this phase is one for every
standard $u\in K$. The real homomorphism $\eta$ is continuous and $K$ is
connected; its image is therefore a connected subset of $\Z$, hence zero.
Its natural ultrapower extension vanishes on every limit element too.
This is (11). Equations (12)--(15), with the \emph{original} vector test
$F_\chi$, now prove the asserted all-box lower-degree representation.

For the torsion assertion, let $\Psi$ be the given degree-at-most-$(d-1)$
vector nilsequence equal to $\chi^{\otimes a}$ on $I^k$.
Put its finitely many coordinates into one common product nilmanifold of
that degree. Since $\chi$ is a unit vector, pointwise on $[N]^k$ one has
\[
              \sum_j (\chi^{\otimes a})_j(n)\overline{\Psi_j(n)}=1.
\]
At least one scalar term has a nonzero standard average, with lower bound
at least the reciprocal of the fixed tensor dimension in absolute value.
Use that term as the test on the product of the original $\chi$ geometry
and the auxiliary geometry of $\Psi$. Its top vertical character is
$a\eta$. Lemma~\ref{lem:integral} applies with the full positive cube and
no nontrivial cutoff. The same Haar argument gives
$a\eta(\operatorname{pr}_\chi K)=0$ as a real homomorphism.
As $a\geq1$, (11) follows. The same reconstruction (12)--(15) uses the
original $F_\chi$, not its tensor power, and proves the required conclusion.

In the symbol formulation of [GTZ, E.13],
$a[\chi]=0$ means exactly that the tensor power is equivalent to the
constant nilcharacter. By the definition of equivalence, its coordinates
are lower-degree nilsequences on the specified box. This is precisely the
$\Psi$ premise just used, with a finite common representation.
No torsion-free conclusion in degree one is asserted; the time and residue
torus factors explain the requirement $d\geq2$.
\end{proof}

\section{Boundaries of this proposal}
This proof supplies the finite-parameter factorization actually needed by
E.11/E.13 and their all-symmetric-box reconstruction. It does not assert
the false horizontal Taylor closeness printed in D.5. Gaussian's separate
one-variable D.6 argument must retain its denominator correction in the
degree-rank horizontal quotient.
The inverse-affine Lemma~\ref{lem:affine} also validates the large-offset
coefficient transition in the one-variable weak-factorization source.
That is a related use requiring an explicit link to this frozen version;
the earlier DF002 bytes and earlier decisions are preserved.
The Section 13 application still needs an independent decision binding its
exact E.11/E.13 objects and domains to this proposal, and the complete
inverse theorem still needs its other separately assigned components.

\begin{thebibliography}{9}
\bibitem{GT14}[GT14] B. Green and T. Tao, On the quantitative distribution
of polynomial nilsequences---erratum, Ann. of Math. 179 (2014), 1175--1183.
DOI: 10.4007/annals.2014.179.3.8.
\bibitem{GT12}[GT12] B. Green and T. Tao, The quantitative behaviour of
polynomial orbits on nilmanifolds, Ann. of Math. 175 (2012), 465--540.
DOI: 10.4007/annals.2012.175.2.2.
\bibitem{GT10}[GT10] B. Green and T. Tao, Linear equations in primes,
Ann. of Math. 171 (2010), 1753--1850. DOI: 10.4007/annals.2010.171.1753.
\bibitem{GTZ}[GTZ] B. Green, T. Tao and T. Ziegler, An inverse theorem for
the Gowers $U^{s+1}[N]$-norm, Ann. of Math. 176 (2012), 1231--1372.
DOI: 10.4007/annals.2012.176.2.11.
\end{thebibliography}
\end{document}
